\section{PC Interface and Sound Output}
\label{sec:pc_interface}

\subsection{UART Frame Format}

The FPGA streams compact ASCII frames to the PC through RS232/UART. Two frame types are used. The $H^2$ frame is mainly useful for debugging and LUT collection:

\begin{lstlisting}
H2,75,H75S1,H75S2,H75S3,45,H45S1,H45S2,H45S3
\end{lstlisting}

The final runtime piano interface uses the key frame:

\begin{lstlisting}
KEY,75,KEY75,PRESS75,45,KEY45,PRESS45
\end{lstlisting}

\code{KEY75} and \code{KEY45} are two-digit global key indices, from 00 to 14. \code{PRESS75} and \code{PRESS45} are debounced active-high button states.

\subsection{Role of the PC}

The PC does not perform magnetic localization in the final system. Its responsibilities are:

\begin{itemize}
    \item Receive the FPGA key/press frame.
    \item Map global key indices to note names.
    \item Display a live piano-key GUI.
    \item Play or stop the corresponding sound when the FPGA press bit changes.
\end{itemize}

This division keeps the project focused on FPGA computation. The Python script is an output interface, not a classifier.

\subsection{Piano GUI and Sound}

The final Python monitor maps global key indices to musical notes. The command used for the 15-key layout is:

\begin{lstlisting}
python3 src_python/piano_key_monitor.py COM5 \
  --notes F3,G3,A3,B3,C4,D4,E4,F4,G4,A4,B4,C5,D5,E5,F5 \
  --buffer-size 64 --poll-ms 2 --volume 0.6
\end{lstlisting}

The monitor starts a note when the corresponding press bit becomes 1 and stops it when the press bit returns to 0. With the pygame audio backend, the sound is played as a held note with low-latency start and exact stop on release.

\placeholderfigure{Insert a screenshot of the Python piano GUI showing highlighted keys.}{PC piano GUI receiving FPGA-computed key states.}{fig:piano_gui}{figures/piano_gui.png}
